\documentclass[a4paper]{article}

\newcommand{\docTitle}{Robotik}

\usepackage{datetime2}  % for dates
\usepackage[utf8]{inputenc} % UTF-8 input encoding
\usepackage[T1]{fontenc} % fixes \textquotedbl with pdflatex/OT1
\usepackage{textcomp}   % additional text symbols for typewriter/upquote
\usepackage{xcolor}     % for gray color
\usepackage{listings}
\usepackage{marginnote} % for horizontal margin notes
\usepackage{parskip}
\usepackage{amsmath,amssymb,mathtools}
\usepackage{everypage}
\usepackage{fancyhdr}
\usepackage{amsfonts}   % Mengen and natural numbers
\usepackage{enumerate}  % Custom Enumerations
\usepackage[inline]{enumitem}   % List spacing and labels
\usepackage{multicol}   % Multiple Columns
\usepackage{tabularx}

\usepackage{hyperref}   % Links
\hypersetup{
    pdftitle={\docTitle},
    pdfauthor={Moritz Köhler},
    pdfkeywords={DHBW, Hochschule},
    colorlinks=true,
    linkcolor=black,
    filecolor=magenta,
    urlcolor=cyan
}

% -------
% Images
% -------

\usepackage{import}
\usepackage{xifthen}
\usepackage{pdfpages}
\usepackage{tikz}
\usetikzlibrary{arrows.meta, positioning}

\newcommand{\incfig}[1]{%
    \def\svgwidth{\columnwidth}
    \import{./fig/}{#1.pdf_tex}
}

% ---------------------------------
% Vertical Side note (Bottom Left)
% ---------------------------------

\newcommand{\verticaldocinfo}{%
  \rotatebox{90}{\tiny\color{gray}\texttt{\DTMtoday\_\docTitle\_\jobname}}%
}

\AddEverypageHook{
  \put(-40pt,-650pt){\verticaldocinfo}%
}

% ---------------
% Math Shortcuts
% ---------------

% Number sets
\newcommand{\R}{\mathbb{R}}
\newcommand{\N}{\mathbb{N}}
\newcommand{\Z}{\mathbb{Z}}
\newcommand{\Q}{\mathbb{Q}}
\newcommand{\C}{\mathbb{C}}

% Vectors and matrices
\newcommand{\vect}[1]{\mathbf{#1}}
\newcommand{\mat}[1]{\mathbf{#1}}

% Derivatives
\newcommand{\dd}{\mathrm{d}}
\newcommand{\dpartial}[2]{\frac{\partial #1}{\partial #2}}
\newcommand{\ddt}{\frac{\dd}{\dd t}}

% Norms and absolute values
\newcommand{\norm}[1]{\left\lVert #1 \right\rVert}
\newcommand{\abs}[1]{\left\lvert #1 \right\rvert}

% Probability & Statistics
\newcommand{\E}{\mathbb{E}}
\newcommand{\Var}{\mathrm{Var}}
\newcommand{\Prob}{\mathbb{P}}

% Fields and Algebras
\newcommand{\K}{\mathbb{K}}
\newcommand{\F}{\mathbb{F}}

% Set Operations
\newcommand{\compl}[1]{#1^{\mathrm{c}}}

% Matrix and Vector Operations
\newcommand{\T}{^{\top}}
\newcommand{\inv}{^{-1}}
\newcommand{\conj}[1]{\overline{#1}}

% -------------------
% Main lecture macro 
% -------------------

\newcommand{\lecture}[1]{%
  \protect\marginnote{\protect\tiny\protect\color{gray}#1}
}

% -----------------------------
% Command for box with heading
% -----------------------------

\fboxsep2mm
\newcommand{\know}[2]{
    \noindent\fbox{
        \parbox{\textwidth-21pt}{

            \textbf{#1:}
            \vspace{0.125cm}

            #2
        }
    }
}

% -------------------
% Listings Languages
% -------------------

% Shared defaults for code listings
\lstset{
    basicstyle=\ttfamily\small,
    keywordstyle=\color{blue}\bfseries,
    commentstyle=\color{gray},
    stringstyle=\color{teal},
    numberstyle=\tiny\color{gray},
    numbers=left,
    stepnumber=1,
    numbersep=8pt,
    showstringspaces=false,
    frame=single,
    breaklines=true,
    tabsize=4,
    keepspaces=false,
    columns=flexible,
    upquote=true,
    extendedchars=false % disable extended char handling to avoid UTF-8 conflicts
}

% SQL syntax highlighting
\lstdefinelanguage{SQL}{
    morekeywords={SELECT,FROM,WHERE,INSERT,INTO,VALUES,UPDATE,SET,DELETE,CREATE,TABLE,PRIMARY,KEY,NOT,NULL,AND,OR,JOIN,LEFT,RIGHT,INNER,ON,AS,ORDER,BY,GROUP,HAVING,LIMIT},
    sensitive=false,
    morecomment=[l]{--},
    morestring=[b]'
}

% JSON syntax highlighting
\lstdefinelanguage{json}{
    morestring=[b]",
    stringstyle=\color{teal},
    comment=[l]{//},
    morecomment=[s]{/*}{*/},
        morekeywords={true,false,null},
    literate=
     *{0}{{{\color{blue}0}}}{1}
      {1}{{{\color{blue}1}}}{1}
      {2}{{{\color{blue}2}}}{1}
      {3}{{{\color{blue}3}}}{1}
      {4}{{{\color{blue}4}}}{1}
      {5}{{{\color{blue}5}}}{1}
      {6}{{{\color{blue}6}}}{1}
      {7}{{{\color{blue}7}}}{1}
      {8}{{{\color{blue}8}}}{1}
      {9}{{{\color{blue}9}}}{1}
            {:}{{{\color{black}:}}}{1}
            {,}{{{\color{black},}}}{1},
    showstringspaces=false
}

% Language specific styles
\lstdefinestyle{javastyle}{
    language=Java,
    basicstyle=\ttfamily\small,
    keywordstyle=\color{blue}\bfseries,
    commentstyle=\color{gray},
    stringstyle=\color{teal},
    numberstyle=\tiny\color{gray},
    numbers=left,
    stepnumber=1,
    numbersep=8pt,
    showstringspaces=false,
    frame=single,
    breaklines=true,
    tabsize=4,
    keepspaces=true,
    columns=fullflexible,
    upquote=true
}

\lstdefinestyle{sqlstyle}{
    language=SQL,
    basicstyle=\ttfamily\small,
    keywordstyle=\color{blue}\bfseries,
    commentstyle=\color{gray},
    stringstyle=\color{teal},
    numberstyle=\tiny\color{gray},
    numbers=left,
    stepnumber=1,
    numbersep=8pt,
    showstringspaces=false,
    frame=single,
    breaklines=true,
    tabsize=4,
    keepspaces=true,
    columns=fullflexible,
    upquote=true
}

\lstdefinestyle{pythonstyle}{
    language=Python,
    basicstyle=\ttfamily\small,
    keywordstyle=\color{blue}\bfseries,
    commentstyle=\color{gray},
    stringstyle=\color{teal},
    numberstyle=\tiny\color{gray},
    numbers=left,
    stepnumber=1,
    numbersep=8pt,
    showstringspaces=false,
    frame=single,
    breaklines=true,
    tabsize=4,
    keepspaces=true,
    columns=fullflexible,
    upquote=true
}

\lstdefinestyle{cstyle}{
    language=C,
    basicstyle=\ttfamily\small,
    keywordstyle=\color{blue}\bfseries,
    commentstyle=\color{gray},
    stringstyle=\color{teal},
    numberstyle=\tiny\color{gray},
    numbers=left,
    stepnumber=1,
    numbersep=8pt,
    showstringspaces=false,
    frame=single,
    breaklines=true,
    tabsize=4,
    keepspaces=true,
    columns=fullflexible,
    upquote=true
}

\lstdefinestyle{bashstyle}{
    language=bash,
    basicstyle=\ttfamily\small,
    keywordstyle=\color{blue}\bfseries,
    commentstyle=\color{gray},
    stringstyle=\color{teal},
    numberstyle=\tiny\color{gray},
    numbers=left,
    stepnumber=1,
    numbersep=8pt,
    showstringspaces=false,
    frame=single,
    breaklines=true,
    tabsize=4,
    keepspaces=true,
    columns=fullflexible,
    upquote=true
}

\lstdefinestyle{jsonstyle}{
    language=json,
    basicstyle=\ttfamily\small,
    keywordstyle=\color{blue}\bfseries,
    commentstyle=\color{gray},
    stringstyle=\color{teal},
    numberstyle=\tiny\color{gray},
    numbers=left,
    stepnumber=1,
    numbersep=8pt,
    showstringspaces=false,
    frame=single,
    breaklines=true,
    tabsize=4,
    keepspaces=true,
    columns=fullflexible,
    upquote=true
}

% Default listing style
\lstset{language={}}

\title{\docTitle}
\author{Moritz}
\date{\today}

\begin{document}
\maketitle
\tableofcontents

\lecture{2026-10-19}

\section{Matlab}

\subsection{Plots}

\begin{lstlisting}
% Test-Daten generieren
time = [0:0.01:120];
f = 0.1;
omega = 2 * pi * f; %Kreisfrequenz rad/s
sig1 = sin(omega * time);
sig2 = cos(omega * time);

figure
plot(time, sig1, 'r', 'LineWidth', 2)
hold on
plot(time, sig2, 'b--', 'LineWidth', 2)
ylim([-1.2, 1.2])
ylabel('Signal [-]')
xlabel('Zeit t [s]')
title('Mein erstes Signal')
grid on
legend
\end{lstlisting}

Mit Subplot und Axis Linking: 

\begin{lstlisting}
figure
title('Mein erstes Signal')
ai(1) = subplot(2, 1, 1);
plot(time, sig1, 'r');
hold on
ylim([-1.2, 1.2])
ylabel('Signal [-]')
xlabel('Zeit t [s]')
grid on
legend

ai(2) = subplot(2, 1, 2);
plot(time, sig2, 'b--');
ylim([-1.2, 1.2])
ylabel('Signal [-]')
xlabel('Zeit t [s]')
grid on
legend

linkaxes(ai, 'x')
\end{lstlisting}

Solver soll \textit{Fixed Step} mit \textit{0.01} und als Solver zuerst \textit{ode1 (Euler)}

\subsection{Simulink einbinden}

\begin{lstlisting}
%% Simulink-Aufruf

x0 = 0;
out = sim("firstSimulink.slx")

figure('Name', 'Simulationsergebnisse')
plot(out.tout, out.results, 'r-')
\end{lstlisting}

Es gibt den Library Browser und wir können die simulationsgeschwindigkeit künstlich auf Echtzeit reduzieren.

\section{Planung, Vorsteuerung und Regelung}

\begin{minipage}[t]{\textwidth}
    \vspace{0pt}
    \centering
    \includegraphics[
        page=113,
        width=\textwidth,
        trim=0pt 50pt 0pt 180pt,
        clip
    ]{doc/Robotik_2026.pdf}
\end{minipage}

Wer diese Grafik versteht, hat Regelung verstanden. Diese hat zwei Freiheitsgrade und kann auf störgrößen reagieren.

Wir kennen unsere Wunschbewegung und die Vorsteuerung macht aus der Bewegung unsere Ansteuerung.

Die Regelung kompensiert Störgrößen und Modellfehler. Es Vergleicht die Wunschbewegung mit der Tatsächlichen Bewegung.

\subsection{Trajektorienplanung}

Es gibt $C^n$-stetige Trajektorien. Das $n$ steht für die Anzahl an Ableitungen die stetig sind.

$PTn$-Filter als Online Trajektorienplaner.

\subsection{Vorsteuerung}

Feedforward Control (Keine Regelschleife)

Bildet keine Modellfehler oder Störgrößen ab.

Die Vorsteuerung nimmt keinen einfluss auf das Stabilitätsverhalten (schwingt es auf oder nicht).

\subsection{Regelung}

Feedback Control (Regelungsschleife)

Braucht Sensorik und verändert das Stabilitätsverhalten.

Mögliche Regler:

\begin{itemize}
    \item Empirisch (z.B. durch händisches Einstellen der PID-Gain)
    \item Modellbasiert, z.B. durch Polvorgabe für ein gegebenes (lineares) Streckenmodell
\end{itemize}

PID (Proportional, Integral, Differential)

Das Dargestellter ist eine Kaskade an Regler.

\begin{minipage}[t]{\textwidth}
    \vspace{0pt}
    \centering
    \includegraphics[
        page=124,
        width=\textwidth,
        trim=0pt 50pt 0pt 55pt,
        clip
    ]{doc/Robotik_2026.pdf}
\end{minipage}

\subsection{Vorsteuerung und Regelung in Gelenkkoordination}

Wir machen unsere Vorsteuerung mit "Computed Torque"-Ansatz.

$$\tau = u = Mv + C\dot q + f + g$$

\end{document}